
\chapter{Preliminaries}

\section{Notations}
\begin{itemize}
    \item $M$: Given an $n \times n$ input matrix $A$, our greedy algorithm seeks to identify a sequence of operations that transforms $A$ into a permutation matrix.
    \item $d$: The circuit depth achieved when matrix $A$ is transformed into a permutation matrix.
    \item $M'$: The matrix cost function is derived from the inverse of matrix $A$.
    \item SIZE: Since we are specifically considering square block matrices, the matrix size is determined by retrieving the row or column length using the built-in `length()` function.
    \item $L_{r}, L_{c}$: Row and column operation sequences are stored in separate layer lists for each iteration.
    \item $Ls_{r}, Ls_{c}$: We maintain row-layer and column-layer lists to store the specific operations assigned to each discrete circuit layer.
    \item $L_{row}, L_{col}$: A list of remaining operations is updated after each step to exclude those already applied to the matrix.
    \item $L_{row}^{cost}, L_{col}^{cost}$: For every operation in $L_{row}$ and $L_{col}$, the algorithm calculates and stores the resulting matrix cost function in a lookup list.
    \item $row_{op}, col_{op}$: The collection of applied row and column operations is stored to construct the final sequence list and to verify the correctness of the transformation.
    \item $p$: A control parameter is used to select the matrix cost function type. Supported functions include linear sum (1), square (2), cubic (3), quartic (4), and logarithmic (- 1), which can be selected via their respective numerical identifiers.
    \item over\_depth: An acceptance flag is utilized to evaluate the current circuit depth; it returns 'True' if the depth $d$ exceeds the current minimum depth, and 'False' otherwise.
    \item close\_permu: A termination flag monitors whether the matrix is nearing a permutation state. If the current matrix can be transformed into a permutation matrix within a single layer, the flag is set to True; otherwise, it remains False.
    \item minm: The algorithm identifies the current minimum cost across all candidates and selects the operation corresponding to this optimal value.
    \item origin, inverse: At the start of the process, a copy of the input matrix is retained for final verification. Simultaneously, the matrix inverse is calculated to facilitate the evaluation of the cost function.
    \item $visi_{row}, visi_{col}$: A visited list tracks the selected row and column indices; once a position is chosen, it is excluded from future iterations to prevent redundant operations.
    \item select\_list: The selection list stores the set of candidate operations available for the current optimization step.
    \item $H_{r}, H_{c}$: The matrix cost function is evaluated from both a row perspective and a column perspective to determine the most efficient operation.
    \item $B_{row}, B_{col}$: The algorithm maintains a collection of row and column operations for both the Row/Column\-wise and Parallel greedy approaches; it identifies the available operators and selects the one with the lowest computational cost.
    \item $best_{row}^{cost}, best_{col}^{cost}$: The algorithm maintains a collection of cost values for each potential operation in $B_{row}$ and $B_{col}$ to facilitate the selection of the optimal operator.
    \item stuck\_counter: A counter is used to record the number of consecutive stalls (stuck iterations) to detect when the algorithm has reached a local optimum.
    \item max\_stuck\_iteration: A stagnation threshold is defined to limit the number of consecutive stalls; by default, this value is set to 3.
    \item escapeCandidates: The algorithm also maintains a collection of non-optimal operators, representing candidates whose cost values exceeded the current minimum.
    \item best\_escape: To resolve local minima, the algorithm identifies the optimal candidate from the escapeCandidates list—after sorting it in ascending order—and adds it to the select\_list to resume the search.
    \item occur\_time: The acceptance counter records the total number of iterations in which an operation was successfully selected and applied.
    \item minm\_size: A threshold condition is implemented to ensure that the circuit depth (total number of layers) does not exceed the predefined limit.
    \item reduce: A copy of the initial matrix is retained to verify that the sequence of row and column operations ($row_{op}$ and $col_{op}$) successfully transforms the origin matrix into the target representation.
    \item size\_minm: The algorithm records the final circuit depth to document the efficiency of the optimized sequence.
    \item per: The algorithm records the final matrix state to generate a permutation list, which defines the final mapping of the quantum registers.
    \item seq: The algorithm maintains a comprehensive log of all applied row and column operations ($row_{op}$ and $col_{op}$), which constitutes the final quantum circuit sequence.
    \item layers: The algorithm maintains a record of the sets $Ls_{c}$ and $Ls_{r}$, which contain the column and row operations assigned to each parallel layer.
    \item correct: A validation flag is employed to confirm that the combined operation sequence and layer assignments successfully transform the original matrix into a permutation matrix.
    \item greedy: An algorithmic selection parameter allows the user to specify which greedy strategy (e.g., Row, Column, Row\_or\_Col or Parallel) should be employed during the synthesis process.
\end{itemize}

\section{Quantum Computation}
